
%% =========================================================
%% Chapter: Symmetries and String Order
%% Develops the symmetry-theoretic consequences of the MPS Fundamental
%% Theorem: virtual gauges, projective bond representations, and
%% string-order criteria for injective tensors.
%% =========================================================

\chapter{Symmetries and String Order}\label{ch:symmetry}

This chapter collects the symmetry-theoretic consequences of the MPS
Fundamental Theorem~\cite{PerezGarcia2007Matrix}, following the symmetry
analysis of~\cite{PerezGarcia2008StringOrder}. For injective tensors,
physical on-site symmetries determine virtual gauges, those gauges define
projective representations on the bond space, and the same virtual data
controls string-order phenomena.

\section{Physical Symmetries Give Virtual Gauges}

For periodic tensors, the physical-index rotation track is developed in
Chapter~\ref{ch:periodic_ft_apps} (which follows this chapter):
Definition~\ref{def:rotate_physical}
and Theorems~\ref{thm:transfer_map_rotate_physical},
\ref{thm:left_canonical_rotate_physical},
\ref{thm:irreducible_tensor_rotate_physical},
\ref{thm:periodic_rotate_physical}, \ref{thm:same_mpv_rotate_physical},
\ref{thm:rotate_physical_blocks}, and
\ref{thm:irreducible_form_rotate_physical}. We use these results here
when deriving symmetry consequences; their proofs are given in
Chapter~\ref{ch:periodic_ft_apps}.

\begin{corollary}[Periodic symmetry yields gauge equivalence up to a root of unity]%
    \label{cor:symmetry_periodic_assembly}
    \lean{MPSTensor.zGaugeEquiv_of_isIrreducibleForm_sameMPV_rotatePhysical}
    \leanok
    \uses{def:rotate_physical, def:irreducible_form, def:same_mpv, thm:periodic_ft}
    Let $A$ be in irreducible form~II and let $u$ be a matrix such that
    $B := A_u$ is also in irreducible form~II and
    $\mathcal{V}(A) = \mathcal{V}(B)$.
    Then there exists $m > 0$ such that $A$ and~$B$ are
    gauge equivalent up to an $m$-th root of unity.
\end{corollary}

\begin{proof}\leanok
    \uses{thm:periodic_ft}
    Direct application of the periodic equal-case fundamental theorem
    (Theorem~\ref{thm:periodic_ft}) to~$A$ and~$B$.
\end{proof}

\section{Virtual Representation Theorem}\label{sec:injective_symmetry_cocycle}

With a full group of on-site symmetries (rather than just a single
unitary), the resulting gauges determine a projective representation
on the bond space.

\begin{definition}[Twisted tensor]\label{def:twisted_tensor}
    \lean{MPSTensor.twistedTensor}
    \leanok
    \uses{def:mps_tensor}
    Given a group representation $U : G \to \GL_d(\C)$ on the physical
    index, the \emph{$g$-twisted tensor} is defined, for each
    $i \in \{0,\ldots,d{-}1\}$, by
    \[
        \widetilde{A}_g^i :=  \sum_{j=0}^{d-1} U(g)_{ij}\, A^j
    \]
\end{definition}

\begin{definition}[On-site symmetry]\label{def:on_site_symmetric}
    \lean{MPSTensor.IsOnSiteSymmetric}
    \leanok
    \uses{def:twisted_tensor, def:same_mpv}
    A tensor $A$ is \emph{on-site symmetric} under a group
    representation $U : G \to \GL_d(\C)$ if $\mc{V}(A) = \mc{V}(\widetilde{A}_g)$
    for every $g \in G$.
\end{definition}

\begin{lemma}[Functoriality of twisting]\label{lem:twisted_tensor_mul}
    \lean{MPSTensor.twistedTensor_mul}
    \leanok
    \uses{def:twisted_tensor}
    The twisted tensor satisfies the composition law
    \[
        \widetilde{A}_{gh} =  \widetilde{(\widetilde{A}_h)}_g,
    \]
    i.e.\ twisting by $gh$ is the same as first twisting by~$h$
    and then by~$g$.
\end{lemma}

\begin{proof}\leanok\uses{def:twisted_tensor}
    Expanding both sides using $U(gh) = U(g)\, U(h)$, then swapping
    the order of summation.
\end{proof}

\begin{lemma}[Identity twist]\label{lem:twisted_tensor_one}
    \lean{MPSTensor.twistedTensor_one}
    \leanok
    \uses{def:twisted_tensor}
    Twisting by the identity group element is trivial:
    $\widetilde{A}_1 = A$.
\end{lemma}

\begin{proof}\leanok\uses{def:twisted_tensor}
    Since $U(1) = I_d$, each component reduces to
    $\sum_j \delta_{ij}\, A^j = A^i$.
\end{proof}

\begin{lemma}[Symmetric twists are gauge equivalent]%
    \label{lem:gauge_equiv_twisted}
    \lean{MPSTensor.gaugeEquiv_twistedTensor_of_injective}
    \leanok
    \uses{def:injective, def:on_site_symmetric}
    If $A$ is injective and on-site symmetric under $U$, then
    for every $g \in G$ the twisted tensor $\widetilde{A}_g$ is
    gauge equivalent to~$A$.
\end{lemma}

\begin{proof}\leanok\uses{def:on_site_symmetric, thm:ft_single}
    On-site symmetry gives
    $\mc{V}(A) = \mc{V}(\widetilde{A}_g)$, and the single-block
    Fundamental Theorem turns this equality of MPV families into gauge
    equivalence.
\end{proof}

\begin{remark}[Fundamental-Theorem route for injective symmetry]%
    \label{rem:injective_symmetry_ft_route}
    \uses{thm:ft_single, lem:gauge_equiv_twisted, lem:gauge_unique_up_to_scalar,
        lem:twisted_tensor_mul, thm:virtual_rep_injective}
    The injective virtual-representation theorem is not proved from string-order
    rigidity. The on-site symmetry condition first identifies $A$ and
    $\widetilde{A}_g$ as the same MPV family; the single-block Fundamental
    Theorem gives the gauge in Lemma~\ref{lem:gauge_equiv_twisted}. Gauge
    uniqueness and the twist composition law then assemble these gauges into
    Theorem~\ref{thm:virtual_rep_injective}. The string-order results later in
    this chapter use this virtual representation as an input.
\end{remark}

\begin{theorem}[Virtual symmetry equation]\label{thm:virtual_symmetry_eq}
    \lean{MPSTensor.virtual_symmetry_eq}
    \leanok
    \uses{def:injective, def:on_site_symmetric, lem:gauge_equiv_twisted}
    If $A$ is injective and on-site symmetric under $U$, then
    for every $g \in G$ there exist an invertible matrix $X(g) \in \GL_D(\C)$
    and a nonzero scalar $\varphi(g) \in \C^{\times}$ (in fact $\varphi = 1$
    in the single-block case) such that
    $\widetilde{A}_g^{\,i} = \varphi(g)\, X(g)\, A^i\, X(g)^{-1}$ for all~$i$.
\end{theorem}

\begin{proof}\leanok\uses{lem:gauge_equiv_twisted}
    The physical action is transferred to the virtual level by the same
    local move that replaces the twisted tensor by a gauge-conjugated one:
    % Use a physical-leg operator insertion for the
    % linear twist \widetilde{A}_g^{\,i} = \sum_j U(g)_{ij}\, A^j rather
    % than the relabelling-only permutation-twist glyph.
    \begin{center}
        \TNLinearTwist{U(g)}{i}
        \qquad \(\leadsto\) \qquad
        \TNGaugeConjugation{X(g)}{i}{X(g)^{-1}}
    \end{center}
    Apply Lemma~\ref{lem:gauge_equiv_twisted} to obtain $X(g)$, then
    set $\varphi(g) = 1$.
\end{proof}

\begin{lemma}[Gauge uniqueness up to scalar]\label{lem:gauge_unique_up_to_scalar}
    \lean{MPSTensor.gauge_unique_up_to_scalar}
    \leanok
    \uses{def:injective, def:gauge_equiv, lem:center_scalar}
    Let $A$ be an injective tensor.
    If $X, Y \in \GL_D(\C)$ both satisfy
    $B^i = X\, A^i\, X^{-1} = Y\, A^i\, Y^{-1}$
    for all~$i$, then there exists a nonzero scalar $u \in \C^{\times}$
    such that $Y = u \cdot X$.
\end{lemma}

\begin{proof}\leanok\uses{def:injective, lem:center_scalar}
    The matrix $Z = Y^{-1} X$ commutes with every $A^i$.
    Since $A$ is injective, $\{A^i\}$ spans $\MN{D}$, so $Z$
    lies in the centre of $\MN{D}$.
    By Lemma~\ref{lem:center_scalar}, $Z$ is a scalar matrix.
    Since $Z$ is invertible, that scalar is nonzero, and
    $Y = u \cdot X$.
\end{proof}

\begin{theorem}[Gauge uniqueness from equality of MPV families]
    \label{thm:gauge_unique_up_to_scalar_sameMPV}
    \lean{MPSTensor.gauge_unique_up_to_scalar_of_sameMPV}
    \leanok
    \uses{def:injective, def:same_mpv}
    Let $A$ be an injective tensor and assume that $A$ and $B$ generate the
    same MPV family.
    If $X, Y \in \GL_D(\C)$ both satisfy
    $B^i = X\, A^i\, X^{-1} = Y\, A^i\, Y^{-1}$ for all $i$, then there is a
    nonzero scalar $u \in \C^{\times}$ such that $Y = u \cdot X$.
\end{theorem}

\begin{proof}
    \leanok
    \uses{lem:gauge_unique_up_to_scalar}
    This is exactly Lemma~\ref{lem:gauge_unique_up_to_scalar}; the extra MPV
    hypothesis specifies the setting in which the two gauges are usually
    produced.
\end{proof}

\begin{definition}[Scalar 2-cochain]\label{def:scalar_cocycle}
    \lean{TNLean.Algebra.ScalarCocycle}
    \leanok
    A \emph{scalar 2-cochain} on a group $G$ is a function
    $\omega : G \times G \to \C^{\times}$.
\end{definition}

\begin{definition}[Projective representation]\label{def:projective_rep}
    \lean{TNLean.Algebra.ProjectiveRepresentation}
    \leanok
    \uses{def:scalar_cocycle}
    Given a scalar 2-cochain $\omega$, a \emph{projective representation}
    of $G$ on $\C^D$ is a map $\rho : G \to \GL_D(\C)$ satisfying, for all
    $g,h \in G$,
    \[
        \rho(g)\, \rho(h) = \omega(g,h)\, \rho(gh)
    \]
\end{definition}

\begin{lemma}[Associativity forces the cocycle condition]%
    \label{lem:cocycle_of_assoc}
    \lean{TNLean.Algebra.ProjectiveRepresentation.cocycle_of_assoc}
    \leanok
    \uses{def:projective_rep, def:scalar_cocycle}
    If $\rho$ is a projective representation with factor system $\omega$
    and $D > 0$, then $\omega$ satisfies the 2-cocycle identity:
    \[
        \omega(g, h)\, \omega(gh, k)
        = \omega(g, hk)\, \omega(h, k).
    \]
\end{lemma}

\begin{proof}\leanok\uses{def:projective_rep}
    Expanding $\rho(g)\,(\rho(h) \rho(k))$ and
    $(\rho(g) \rho(h)) \rho(k)$ using the projective
    multiplication law and equating via matrix associativity.
    The hypothesis $D > 0$ allows cancellation of the
    invertible matrix factor.
\end{proof}

\begin{theorem}[Virtual representation theorem for injective MPS]%
    \label{thm:virtual_rep_injective}
    \lean{MPSTensor.virtual_rep_of_symmetric_injective}
    \leanok
    \uses{def:injective, def:on_site_symmetric, def:projective_rep,
        def:scalar_cocycle, lem:gauge_unique_up_to_scalar,
        lem:twisted_tensor_mul}
    Let $A$ be an injective MPS tensor that is on-site symmetric
    under a group representation $U : G \to \GL_d(\C)$.
    Then there exist:
    \begin{enumerate}
        \item a scalar 2-cochain
              $\omega : G \times G \to \C^{\times}$, and
        \item a map $\rho : G \to \GL_D(\C)$ satisfying, for all $g,h \in G$,
              \[
                  \rho(g)\, \rho(h) = \omega(g,h)\, \rho(gh)
              \]
    \end{enumerate}
    such that, defining the gauge matrices $X(g) := \rho(g^{-1})$,
    one has, for every $g \in G$ and $i \in \{0,\ldots,d{-}1\}$,
    \[
        \widetilde{A}_g^i =  X(g)\, A^i\, X(g)^{-1}
    \]
    In tensor-network notation, for each fixed~$g$ this is the local gauge
    relation
    \begin{center}
        \TNGaugeConjugation{X(g)}{i}{X(g)^{-1}}
    \end{center}
    with the black node denoting the tensor named in the surrounding
    formula and the two red side nodes acting on the virtual legs.
    In particular, $\rho$ is a projective representation of $G$ on the
    bond space $\C^D$.
\end{theorem}

\begin{proof}
    \leanok
    \uses{def:on_site_symmetric, lem:gauge_equiv_twisted,
        lem:gauge_unique_up_to_scalar, lem:twisted_tensor_mul}
    Since $A$ is injective and on-site symmetric, each twisted tensor
    $\widetilde{A}_g$ is gauge equivalent to~$A$
    (Lemma~\ref{lem:gauge_equiv_twisted}).
    Choose $X(g) \in \GL_D(\C)$ with
    $\widetilde{A}_g^i = X(g)\, A^i\, X(g)^{-1}$ for all~$i$.

    By functoriality (Lemma~\ref{lem:twisted_tensor_mul}), the product
    $X(h)\, X(g)$ also intertwines $A$ with $\widetilde{A}_{gh}$.
    Since $X(gh)$ does the same, gauge uniqueness
    (Lemma~\ref{lem:gauge_unique_up_to_scalar}) gives a nonzero scalar
    $\omega'(h,g)$ with $X(h)\, X(g) = \omega'(h,g)\, X(gh)$.

    Defining $\rho(g) := X(g^{-1})$ converts this to the standard law:
    $\rho(g)\, \rho(h) = \omega(g,h)\, \rho(gh)$
    where $\omega(g,h) := \omega'(h^{-1}, g^{-1})$.
\end{proof}

\begin{corollary}[2-cocycle identity for the factor system]%
    \label{cor:virtual_rep_cocycle}
    \lean{MPSTensor.virtual_rep_of_symmetric_injective_cocycle}
    \leanok
    \uses{thm:virtual_rep_injective, lem:cocycle_of_assoc}
    If $D > 0$, the factor system $\omega$ from
    Theorem~\ref{thm:virtual_rep_injective} satisfies the
    2-cocycle identity for all $g,h,k \in G$:
    \[
        \omega(g, h)\, \omega(gh, k)
        = \omega(g, hk)\, \omega(h, k)
    \]
\end{corollary}

\begin{proof}
    \leanok
    \uses{thm:virtual_rep_injective, lem:cocycle_of_assoc}
    Theorem~\ref{thm:virtual_rep_injective} gives $\rho$ and $\omega$.
    By Lemma~\ref{lem:cocycle_of_assoc}, the cocycle identity follows
    from associativity of matrix multiplication and invertibility of
    the representation matrices.
\end{proof}

\section{Cohomology Classes of Virtual Symmetry Data}

Two cocycles may differ by a coboundary yet represent the same
projective-representation class. The following definitions and results
make this precise and establish the quotient $H^2(G,\C^{\times})$.

\begin{definition}[Coboundary]\label{def:is_coboundary}
    \lean{TNLean.Algebra.ScalarCocycle.IsCoboundary}
    \leanok
    \uses{def:scalar_cocycle}
    A scalar 2-cocycle $\omega$ is a \emph{coboundary} if there exists
    $\varphi \colon G \to \C^{\times}$ such that, for all $g,h \in G$,
    \[
        \omega(g,h)
        = \varphi(g) \varphi(h) \varphi(gh)^{-1}
    \]
\end{definition}

\begin{definition}[Cohomologous cocycles]\label{def:cohomologous_to}
    \lean{TNLean.Algebra.ScalarCocycle.CohomologousTo}
    \leanok
    \uses{def:scalar_cocycle}
    Two cocycles $\omega_1,\omega_2$ are \emph{cohomologous}
    (written $\omega_1 \sim \omega_2$) if there exists
    $\varphi\colon G\to\C^{\times}$ such that, for all $g,h \in G$,
    \[
        \omega_1(g,h)
        = \varphi(g) \varphi(h) \varphi(gh)^{-1} \omega_2(g,h)
    \]
\end{definition}

\begin{theorem}[Cohomologous-to is an equivalence relation]%
    \label{thm:cohomologous_equivalence}
    \lean{TNLean.Algebra.ScalarCocycle.CohomologousTo.equivalence}
    \leanok
    \uses{def:cohomologous_to}
    The relation "cohomologous to" is reflexive, symmetric, and
    transitive on the set of scalar 2-cocycles.
\end{theorem}

\begin{proof}\leanok\uses{def:cohomologous_to}
    Reflexivity: take $\varphi = 1$.
    Symmetry: replace $\varphi$ by $\varphi^{-1}$.
    Transitivity: compose the witnesses $\varphi$ and $\psi$ pointwise.
\end{proof}

\begin{lemma}[Coboundary iff cohomologous to the trivial cocycle]%
    \label{lem:coboundary_iff_cohomologous_one}
    \lean{TNLean.Algebra.ScalarCocycle.isCoboundary_iff_cohomologousTo_one}
    \leanok
    \uses{def:is_coboundary, def:cohomologous_to}
    A cocycle $\omega$ is a coboundary if and only if
    $\omega \sim \mathbf{1}$, where $\mathbf{1}(g,h)=1$ for all $g,h$.
\end{lemma}

\begin{proof}\leanok\uses{def:is_coboundary, def:cohomologous_to}
    Both directions follow by absorbing or introducing the factor
    $\mathbf{1}(g,h)=1$ via $x\cdot 1 = x$.
\end{proof}

\begin{theorem}[Gauge independence of the cocycle class]%
    \label{thm:cocycle_gauge_independence}
    \lean{MPSTensor.cohomologousTo_of_isInjective}
    \leanok
    \uses{thm:virtual_rep_injective, lem:gauge_unique_up_to_scalar,
        def:cohomologous_to}
    If the bond dimension satisfies $D \ge 1$ and two projective
    representations $\rho_1,\rho_2$ (with cocycles
    $\omega_1,\omega_2$) both arise as virtual representations of the
    same injective MPS tensor $A$ under an on-site symmetry~$U$, then
    $\omega_1 \sim \omega_2$.
\end{theorem}

\begin{proof}\leanok
    \uses{thm:virtual_rep_injective, lem:gauge_unique_up_to_scalar,
        def:cohomologous_to}
    For each $g$, gauge uniqueness gives a nonzero scalar $f(g)$ with
    $X_2(g^{-1}) = f(g)\,X_1(g^{-1})$.
    Equivalently, for each fixed $g$ the two local gauges
    \begin{center}
        \TNGaugeConjugation{X_1(g^{-1})}{i}{X_1(g^{-1})^{-1}}
        \qquad and\qquad
        \TNGaugeConjugation{X_2(g^{-1})}{i}{X_2(g^{-1})^{-1}}
    \end{center}
    describe the same twisted tensor, so they differ by a scalar.
    Substituting into the projective law and cancelling $X_1(g\cdot h)$
    yields
    $\omega_2(g,h) = f(g)\,f(h)\,f(gh)^{-1} \omega_1(g,h)$.
\end{proof}

\begin{definition}[Multiplicative 2-cocycle condition]\label{def:is_cocycle}
    \lean{TNLean.Algebra.ScalarCocycle.IsCocycle}
    \leanok
    A scalar 2-cochain $\omega\colon G\times G\to\C^{\times}$ is a
    \emph{2-cocycle} if it satisfies, for all $g,h,k \in G$,
    \[
        \omega(g,h) \omega(gh,k) = \omega(g,hk) \omega(h,k)
    \]
\end{definition}

\begin{definition}[Second cohomology quotient]\label{def:h2_cx}
    \lean{TNLean.Algebra.H2}
    \leanok
    \uses{def:cohomologous_to, def:is_cocycle}
    The second cohomology set is the quotient
    \[
        H^2(G,\C^{\times})
        := Z^2(G,\C^{\times})/{\sim},
    \]
    where $\sim$ is the coboundary-equivalence relation.
\end{definition}

\begin{definition}[Projective equivalence]\label{def:projectively_equivalent}
    \lean{TNLean.Algebra.ProjectivelyEquivalent}
    \leanok
    \uses{def:projective_rep, def:cohomologous_to}
    Two projective representations are \emph{projectively equivalent}
    when their factor systems are cohomologous.
\end{definition}

\begin{theorem}[Projective equivalence and cohomologous cocycles]%
    \label{thm:projrep_equiv_iff_cohomologous}
    \lean{TNLean.Algebra.projRep_equiv_iff_cohomologous}
    \leanok
    \uses{def:projectively_equivalent, def:cohomologous_to}
    Let $\rho_1,\rho_2$ be projective representations with factor systems
    $\omega_1,\omega_2$. Then
    \[
        \rho_1 \sim_{\mathrm{proj}} \rho_2
        \iff
        \omega_1 \sim \omega_2.
    \]
\end{theorem}

\begin{proof}\leanok\uses{def:projectively_equivalent}
    This is immediate from the definition of projective equivalence as
    cohomology-class equality for factor systems.
\end{proof}

\section{Permutation-Based Physical Symmetry}

The linear-combination twist $\widetilde{A}_g^i = \sum_j U(g)_{ij}\, A^j$
is the standard form used in the virtual representation theorem.
The following alternative formulation replaces the matrix action by a
permutation of the physical index; it is a formal convenience for
symmetries that act by reordering basis states rather than by a general
linear map.  % The statements of this section are auxiliary constructions not
% directly sourced from~\cite{PerezGarcia2008StringOrder}.

\begin{definition}[Permutation symmetry data]\label{def:on_site_symmetry_perm}
    \lean{MPSTensor.OnSiteSymmetry}
    \leanok
    An \emph{on-site symmetry datum} for a group $G$ on a physical
    space of dimension $d$ consists of a matrix-valued map
    $u : G \to \MN{d}$ and a physical-index action
    $\sigma : G \to \mathrm{End}(\{0,\ldots,d{-}1\})$.
\end{definition}

\begin{definition}[Permutation-twisted tensor]\label{def:perm_twisted_tensor}
    \lean{MPSTensor.TwistedTensor}
    \leanok
    \uses{def:on_site_symmetry_perm, def:mps_tensor}
    Given on-site symmetry data $(u, \sigma)$, the
    \emph{permutation-twisted tensor} at group element $g$ is
    \[
        (A^{\sigma_g})^i :=  A^{\sigma(g)(i)}.
    \]
    In tensor-network notation the local tensor is unchanged and only the
    physical leg is relabelled:
    \begin{center}
        \TNPermutationTwist{i}{\sigma(g)(i)}
    \end{center}
\end{definition}

\begin{lemma}[Identity permutation twist]\label{lem:perm_twisted_one}
    \lean{MPSTensor.TwistedTensor_one}
    \leanok
    \uses{def:perm_twisted_tensor}
    If $\sigma(1) = \mathrm{id}$, then $(A^{\sigma_1})^i = A^i$ for all~$i$.
\end{lemma}

\begin{proof}\leanok\uses{def:perm_twisted_tensor}
    Immediate from $\sigma(1)(i) = i$ for all~$i$.
\end{proof}

\noindent\emph{Composition order.}
The linear-combination twist (Lemma~\ref{lem:twisted_tensor_mul}) composes as
$\widetilde{(\widetilde{A}_h)}_g = \widetilde{A}_{gh}$, applying $h$ first
then~$g$. The permutation twist below uses the same left-action convention
$\sigma(gh) = \sigma(g) \circ \sigma(h)$, so
$(A^{\sigma_g})^{\sigma_h} = A^{\sigma_{gh}}$; the inner permutation $\sigma(h)$
acts first on the index.

\begin{lemma}[Composition of permutation twists]\label{lem:perm_twisted_mul}
    \lean{MPSTensor.TwistedTensor_mul}
    \leanok
    \uses{def:perm_twisted_tensor}
    If $\sigma(gh) = \sigma(g) \circ \sigma(h)$ for all $g, h \in G$,
    then
    \[
        A^{\sigma_{gh}} =  (A^{\sigma_g})^{\sigma_h}.
    \]
\end{lemma}

\begin{proof}\leanok\uses{def:perm_twisted_tensor}
    Expanding both sides:
    $(A^{\sigma_{gh}})^i = A^{\sigma(g)(\sigma(h)(i))} = ((A^{\sigma_g})^{\sigma_h})^i$.
    Diagrammatically this is just two consecutive relabellings of the same
    local tensor:
    % Label the two arrows by \sigma_h and \sigma_g so
    % the relabelling order matches the composition \sigma(g) \circ \sigma(h).
    \begin{center}
        \TNPermutationTwistLabeled{i}{\sigma(h)(i)}{\sigma_h}
        \qquad
        \TNPermutationTwistLabeled{\sigma(h)(i)}{\sigma(g)(\sigma(h)(i))}{\sigma_g}
    \end{center}
\end{proof}

\section{String Order and Local Symmetry Equivalence}

The string order parameter \cite{PerezGarcia2008StringOrder} detects
whether an on-site unitary $u$ acts as a local symmetry on the state
generated by an injective MPS tensor.

\begin{definition}[Twisted transfer map]\label{def:twisted_transfer_map}
    \lean{MPSTensor.twistedTransferMap}
    \leanok
    \uses{def:mps_tensor}
    Given an MPS tensor $A$ and a physical-index matrix $u$, the
    \emph{twisted transfer map} is
    \[
        \E_u(X) :=  \sum_{n,n'} \langle n'|u|n\rangle
        A_n\, X\, A_{n'}^\dagger.
    \]
    Diagrammatically one inserts the physical operator $u$ on the doubled
    local transfer cell:
    \begin{center}
        \TNTwistedTransfer{u}
    \end{center}
    The upper and lower black nodes denote $A_n$ and $A_{n'}^\dagger$, and the
    red node labelled $u$ couples the physical legs.
\end{definition}

\begin{definition}[String order parameter]\label{def:string_order_param}
    \lean{MPSTensor.stringOrderParam}
    \leanok
    \uses{def:twisted_transfer_map}
    The \emph{string order parameter} for twist $u$ and stationary
    state $\Lambda$ at length $L$ is
    \[
        R_L(u) :=  \tr\!(\Lambda \cdot \E_u^L(\Id)).
    \]
    In tensor-network notation this is the length-$L$ doubled chain with a
    copy of $u$ inserted on each physical rung:
    \begin{center}
        \TNStringOrderParameter{u}{L}
    \end{center}
    The boundary matrix $\Lambda$ is contracted into the virtual boundary, the
    red nodes labelled $u$ sit on the physical rungs, and the right boundary
    is traced.
\end{definition}

\begin{definition}[Boundary string order parameter]%
    \label{def:string_order_boundary_param}
    \lean{MPSTensor.stringOrderBoundaryParam}
    \leanok
    \uses{def:twisted_transfer_map, def:string_order_param}
    For a boundary state $\Lambda$ and boundary matrices $X,Y \in \MN{D}$, the
    \emph{boundary string order parameter} is
    \[
        R_L(u;X,Y) := \tr\!\bigl(\Lambda X \, \E_u^L(Y)\bigr).
    \]
    The ordinary string order parameter is the special case $X=Y=\Id$.
\end{definition}

\begin{definition}[Local symmetry]\label{def:local_symmetry}
    \lean{MPSTensor.IsLocalSymmetry}
    \leanok
    \uses{def:mps_tensor}
    An MPS tensor $A$ has \emph{local symmetry} under a unitary $u$ with
    respect to a boundary state $\Lambda$ if there exist a unitary
    matrix $V$ and a phase $\mu$ with $|\mu| = 1$ such that
    $V^\dagger \Lambda V = \Lambda$ and, for every physical index~$i$,
    \[
        \sum_j u_{ij}\, A^j = \mu\, V\, A^i\, V^\dagger .
    \]
    The unitary $V$ intertwines the physical action of~$u$ with a virtual
    conjugation, and the boundary state $\Lambda$ is invariant under~$V$.
\end{definition}

\begin{definition}[String order]\label{def:has_string_order}
    \lean{MPSTensor.HasStringOrder}
    \leanok
    \uses{def:string_order_boundary_param}
    String order \emph{exists} for $u$ with respect to a boundary
    state $\Lambda$ if there exist boundary matrices $X, Y$ and a
    positive constant $c > 0$ such that
    $c \le \|R_L(u; X, Y)\|$ for all $L$, where
    $R_L(u; X, Y)$ is the boundary-refined string-order parameter.
    Informally, this means the scalar string-order parameter
    $R_L(u)$ does not decay to zero.
\end{definition}

\begin{lemma}[String order prevents decay to zero]
    \label{lem:not_tendsto_zero_of_hasStringOrder}
    \lean{MPSTensor.not_tendsto_zero_of_hasStringOrder}
    \leanok
    \uses{def:has_string_order}
    If $A$ has string order for $u$ with respect to a boundary state
    $\Lambda$, then for some boundary matrices $X,Y$ the sequence
    $R_L(u;X,Y)$ does not tend to $0$ as $L \to \infty$.
\end{lemma}

\begin{proof}
    \leanok
    \uses{def:has_string_order}
    A uniform lower bound $c \le \|R_L(u;X,Y)\|$ is incompatible with
    convergence to $0$.
\end{proof}

\begin{definition}[Local physical intertwining]\label{def:condC1}
    \lean{MPSTensor.CondC1}
    \leanok
    \uses{def:mps_tensor}
    The local physical intertwining relation is
    $\sum_j u_{ij}\, A^j = V\, A^i\, V^\dagger$.
    Locally, the physical action of $u$ can therefore be pushed to the virtual
    legs:
    \begin{center}
        \TNCondCOne{u}{V}{i}
    \end{center}
\end{definition}

\begin{definition}[Transfer-map covariance]\label{def:condC2}
    \lean{MPSTensor.CondC2}
    \leanok
    \uses{def:transfer_map}
    Transfer-map covariance is the identity that the transfer map commutes with virtual
    conjugation: $\E(V X V^\dagger) = V\, \E(X)\, V^\dagger$.
    In doubled-layer notation, the virtual operators slide through the local
    transfer cell:
    \begin{center}
        \TNCondCTwo{V}
    \end{center}
\end{definition}

\begin{definition}[Commutation in the doubled transfer picture]\label{def:condC3}
    \lean{MPSTensor.CondC3}
    \leanok
    \uses{def:twisted_transfer_map}
    In the transfer-map language, the commutation of the doubled transfer matrix
    $E = \sum_i A_i \otimes \bar A_i$ with $V \otimes \bar V$ is
    \[
        V\, \E_1(X)\, V^\dagger = \E_1(V X V^\dagger) .
    \]
\end{definition}

\begin{theorem}[Transfer-map covariance and doubled commutation are equivalent]%
    \label{thm:condC2_iff_condC3}
    \lean{MPSTensor.condC2_iff_condC3}
    \leanok
    \uses{def:condC2, def:condC3}
    Transfer-map covariance and commutation in the doubled transfer
    picture are equivalent for any matrix $V$.
\end{theorem}

\begin{proof}
    \leanok
    Both sides express the same local picture with opposite orientation:
    \begin{center}
        \TNCondCTwo{V}
    \end{center}
    The two formulas are literally the same identity read from opposite sides.
\end{proof}

\begin{lemma}[Unitary Kraus mixing]\label{lem:unitary_kraus_mixing}
    \lean{MPSTensor.unitary_kraus_mixing}
    \leanok
    If $u$ is unitary, then replacing Kraus operators $\{A_i\}$ by
    their unitary mixtures $\{\sum_j u_{ij} A_j\}$ preserves the
    channel: $\sum_i (\sum_j u_{ij} A_j)\, Y\, (\sum_j u_{ij} A_j)^\dagger
        = \sum_i A_i\, Y\, A_i^\dagger$.
\end{lemma}

\begin{proof}
    \leanok
    \uses{thm:kraus_same_map_of_unitary_combination}
    This is exactly the usual invariance of a Kraus decomposition under a
    unitary change of Kraus operators; see Theorem~\ref{thm:kraus_same_map_of_unitary_combination}.
\end{proof}

\begin{theorem}[Local physical intertwining implies transfer-map covariance]%
    \label{thm:condC1_imp_condC2}
    \lean{MPSTensor.condC1_imp_condC2}
    \leanok
    \uses{def:condC1, def:condC2, lem:unitary_kraus_mixing}
    If $V$ and $u$ are unitary and the local physical intertwining
    relation holds, then the transfer map is covariant.
\end{theorem}

\begin{proof}
    \leanok
    \uses{lem:unitary_kraus_mixing}
    Starting from
    \begin{center}
        \TNCondCOne{u}{V}{i}
    \end{center}
    one inserts $V^\dagger V = \Id$ so that each local transfer cell is
    written in terms of conjugated Kraus operators, then replaces those
    operators using the intertwining relation. After that, the physical
    $u$-boxes disappear by the unitary Kraus-mixing identity.
\end{proof}

\begin{theorem}[Transfer-map covariance implies local physical intertwining]%
    \label{thm:condC2_imp_condC1_of_injective}
    \lean{MPSTensor.condC2_imp_condC1_of_injective}
    \leanok
    \uses{def:condC2, def:condC1, def:injective}
    For an injective MPS, if $V$ is unitary and the transfer map is
    covariant, then there exists a unitary $u$ satisfying the local
    physical intertwining relation.
\end{theorem}

\begin{proof}
    \leanok
    \uses{thm:condC2_iff_condC3, thm:condC1_imp_condC2}
    Set $B^i := V\, A^i\, V^\dagger$. The transfer-map covariance relation
    \begin{center}
        \TNCondCTwo{V}
    \end{center}
    implies that the two Kraus families $B$ and $A$ define the same
    transfer map. The missing step --- supplied by Kraus freedom for
    equal channels --- converts equality of transfer maps $\E_B = \E_A$
    into a unitary Kraus mixing matrix $u$ with
    $B^i = \sum_j u_{ij} A^j$.
\end{proof}

\begin{definition}[Twisted companion tensor]\label{def:twisted_mixed_companion}
    \lean{MPSTensor.twistedMixedCompanion}
    \leanok
    \uses{def:mps_tensor}
    For an MPS tensor $A$ and a physical-index matrix $u$, define the
    companion tensor $B_u$ by
    \[
        B_u^n := \sum_{n'=0}^{d-1} \overline{u_{n'n}}\, A^{n'}.
    \]
\end{definition}

\begin{lemma}[Twisted transfer as a mixed transfer operator]%
    \label{lem:twisted_transfer_eq_mixed_transfer}
    \lean{MPSTensor.twistedTransferMap_eq_mixedTransfer}
    \leanok
    \uses{def:twisted_transfer_map, def:twisted_mixed_companion, def:mixed_transfer}
    The twisted transfer map of $A$ is the mixed transfer operator of the pair
    $(A,B_u)$:
    \[
        \E_u = F_{A,B_u}.
    \]
\end{lemma}

\begin{proof}
    \leanok
    Expand both definitions and rearrange the double sum.
\end{proof}

\begin{theorem}[Twisted-transfer eigenvalues have modulus at most one]%
    \label{thm:twisted_transfer_spectral_radius_le_one}
    \lean{MPSTensor.twistedTransfer_spectralRadius_le_one}
    \leanok
    \uses{def:twisted_transfer_map, def:injective}
    Let $A$ be an injective MPS tensor with $\E_A(\Id)=\Id$, and let $u$
    satisfy $u u^\dagger = \Id$.
    If
    \[
        \E_u(X) = \lambda X
    \]
    for some nonzero $X \in \MN{D}$, then $|\lambda| \le 1$.
\end{theorem}

\begin{proof}
    \leanok
    \uses{lem:unitary_kraus_mixing, lem:twisted_transfer_eq_mixed_transfer,
        thm:left_canonical_gauge, thm:eigenvalue_bound}
    By Lemma~\ref{lem:twisted_transfer_eq_mixed_transfer}, the twisted transfer
    map is the mixed transfer operator $F_{A,B_u}$.
    Because $B_u$ is obtained from $A$ by unitary mixing of the Kraus operators,
    the transfer maps of $A$ and $B_u$ coincide, so the two adjoint channels
    share a common positive definite fixed point.
    Gauging both families by that fixed point puts them in a common TP gauge.
    Theorem~\ref{thm:eigenvalue_bound} then applies to the mixed transfer
    operator in that gauge, and similarity preserves the eigenvalue $\lambda$.
\end{proof}

\begin{theorem}[Modulus-one twisted eigenvalue implies gauge-phase equivalence]%
    \label{thm:twisted_transfer_modulus_one_gauge_phase}
    \lean{MPSTensor.twistedTransfer_modulus_one_implies_gaugePhase}
    \leanok
    \uses{def:twisted_transfer_map, def:twisted_mixed_companion,
        def:gauge_phase_equiv, def:injective}
    Let $A$ be an injective MPS tensor with $\E_A(\Id)=\Id$, and let $u$
    satisfy $u u^\dagger = \Id$.
    If
    \[
        \E_u(X) = \lambda X
    \]
    for some nonzero $X \in \MN{D}$ with $|\lambda|=1$, then $A$ is
    gauge-phase equivalent to the companion tensor $B_u$.
\end{theorem}

\begin{proof}
    \leanok
    \uses{lem:unitary_kraus_mixing, lem:twisted_transfer_eq_mixed_transfer,
        thm:left_canonical_gauge, thm:modulus_one_gauge_irr_tp}
    By Lemma~\ref{lem:twisted_transfer_eq_mixed_transfer}, the twisted transfer
    map is the mixed transfer operator $F_{A,B_u}$.
    Because $B_u$ is obtained from $A$ by unitary mixing of the Kraus operators,
    the transfer maps of $A$ and $B_u$ coincide, so the two adjoint channels
    share a common positive definite fixed point.
    Gauging both families by that fixed point puts them in a common TP gauge.
    In that gauge the eigenvalue equation becomes a modulus-one peripheral
    eigenvalue for an irreducible mixed transfer operator, so
    Theorem~\ref{thm:modulus_one_gauge_irr_tp} yields gauge-phase equivalence.
    Undoing the common TP gauge gives the claim for $A$ and $B_u$.
\end{proof}

\begin{lemma}[Spectral radius $<1$ forces boundary-string decay]%
    \label{lem:boundary_string_order_tendsto_zero}
    \lean{MPSTensor.stringOrderBoundaryParam_tendsto_zero_of_spectralRadius_lt_one}
    \leanok
    \uses{def:string_order_boundary_param, def:twisted_transfer_map}
    If $\rho(\E_u) < 1$, then for all boundary matrices
    $X,Y,\Lambda \in \MN{D}$ one has
    \[
        R_L(u;X,Y) \longrightarrow 0
        \qquad\text{as } L \to \infty.
    \]
\end{lemma}

\begin{proof}
    \leanok
    In finite dimension, $\rho(\E_u) < 1$ implies $\E_u^L \to 0$ in operator
    norm. Composing with the fixed trace pairing
    $Z \mapsto \tr(\Lambda X Z)$ gives the stated scalar convergence.
\end{proof}

\begin{theorem}[String order gives gauge-phase equivalence with the companion tensor]%
    \label{thm:string_order_implies_twisted_gauge_phase}
    \lean{MPSTensor.gaugePhaseEquiv_twisted_of_hasStringOrder}
    \leanok
    \uses{def:has_string_order, def:twisted_mixed_companion,
        def:gauge_phase_equiv, def:injective}
    Let $A$ be injective, let $\Lambda \in \MN{D}$, and assume
    $u u^\dagger = \Id$ and $\E_A(\Id)=\Id$.
    If string order exists for $u$ with boundary parameter $\Lambda$, then $A$
    is gauge-phase equivalent to the companion tensor $B_u$.
\end{theorem}

\begin{proof}
    \leanok
    \uses{lem:boundary_string_order_tendsto_zero,
        thm:twisted_transfer_spectral_radius_le_one,
        thm:twisted_transfer_modulus_one_gauge_phase}
    If $\rho(\E_u) < 1$, Lemma~\ref{lem:boundary_string_order_tendsto_zero}
    forces every boundary string-order sequence to decay to zero, contradicting
    Definition~\ref{def:has_string_order}.
    Hence $\rho(\E_u) \ge 1$.
    Theorem~\ref{thm:twisted_transfer_spectral_radius_le_one} shows that every
    eigenvalue has modulus at most $1$, so some peripheral eigenvalue has modulus
    exactly $1$.
    Applying Theorem~\ref{thm:twisted_transfer_modulus_one_gauge_phase} then
    gives the gauge-phase equivalence.
\end{proof}

\begin{theorem}[Gauge-phase equivalence with the companion tensor yields a virtual unitary]%
    \label{thm:virtual_unitary_of_twisted_gauge_phase}
    \lean{MPSTensor.virtualUnitary_of_gaugePhaseEquiv_twisted}
    \leanok
    \uses{def:gauge_phase_equiv, def:twisted_mixed_companion, def:injective}
    Let $A$ be injective, assume $u u^\dagger = \Id$ and $\E_A(\Id)=\Id$.
    If $A$ is gauge-phase equivalent to the companion tensor $B_u$, then there
    exist a unitary $V$ and a phase $\mu$ with $|\mu|=1$ such that, for every
    physical index~$i$,
    \[
        \sum_j u_{ij} A^j = \mu\, V A^i V^\dagger
    \]
\end{theorem}

\begin{proof}
    \leanok
    \uses{thm:psd_fp_unique}
    Start from an invertible intertwiner $X$ between $A$ and $B_u$.
    The positive matrix $Q := X X^\dagger$ is a positive fixed point of
    $\E_A$.
    By uniqueness of positive fixed points for an injective normalized tensor,
    $Q$ is a positive scalar multiple of the identity.
    Rescaling $X$ by the square root of that scalar produces a unitary $V$,
    and the original intertwining relation becomes the phased covariance
    equation above, with $|\mu|=1$.
\end{proof}

\begin{theorem}[Virtual unitary gives a twisted-transfer eigenmatrix]%
    \label{thm:twisted_transfer_eigen_of_virtual_unitary}
    \lean{MPSTensor.twistedTransfer_eigen_of_virtualUnitary}
    \leanok
    \uses{def:twisted_transfer_map}
    Assume $\E_A(\Id)=\Id$ and that, for every physical index~$i$,
    \[
        \sum_j u_{ij} A^j = \mu\, V A^i V^\dagger
    \]
    with $V$ unitary.
    Then $V$ is an eigenmatrix of the twisted transfer map:
    \[
        \E_u(V) = \mu V.
    \]
\end{theorem}

\begin{proof}
    \leanok
    Expand $\E_u(V)$, insert the intertwining relation into each local term,
    and use $V^\dagger V = \Id$ together with $\E_A(\Id)=\Id$.
\end{proof}

\begin{theorem}[Boundary-state invariance of the virtual unitary]%
    \label{thm:boundary_state_invariant_of_virtual_unitary}
    \lean{MPSTensor.boundaryState_invariant_of_virtualUnitary}
    \leanok
    \uses{def:injective}
    Let $A$ be injective, assume $u u^\dagger = \Id$, and let $\Lambda$ be a
    positive definite boundary state with $\tr(\Lambda)=1$ and
    $\E_A^\dagger(\Lambda)=\Lambda$.
    If $V$ is a unitary and $\mu$ a phase satisfying, for every physical
    index~$i$,
    \[
        \sum_j u_{ij} A^j = \mu\, V A^i V^\dagger
    \]
    then
    \[
        V^\dagger \Lambda V = \Lambda.
    \]
\end{theorem}

\begin{proof}
    \leanok
    \uses{thm:psd_fp_unique}
    The covariance relation shows that $V^\dagger \Lambda V$ is another positive
    fixed point of the adjoint transfer map.
    Since $\Lambda$ is the normalized positive fixed point, uniqueness forces
    $V^\dagger \Lambda V = \Lambda$.
\end{proof}

\begin{lemma}[Local symmetry implies string order]%
    \label{lem:has_string_order_of_local_symmetry}
    \lean{MPSTensor.hasStringOrder_of_localSymmetry}
    \leanok
    \uses{def:local_symmetry, def:has_string_order}
    Assume $\tr(\Lambda)=1$ and $\E_A(\Id)=\Id$.
    If $u$ is a local symmetry of $A$ with respect to $\Lambda$, then string
    order exists for $u$.
\end{lemma}

\begin{proof}
    \leanok
    \uses{thm:twisted_transfer_eigen_of_virtual_unitary,
        def:string_order_boundary_param}
    Choose the boundary matrices $X = V^\dagger$ and $Y = V$, where $V$ is the
    unitary from the local-symmetry datum.
    By Theorem~\ref{thm:twisted_transfer_eigen_of_virtual_unitary},
    one has $\E_u^L(V)=\mu^L V$.
    Therefore
    \[
        R_L(u;V^\dagger,V)=\mu^L \tr(\Lambda)=\mu^L,
    \]
    whose norm is constantly $1$.
\end{proof}

\begin{theorem}[Local symmetry iff a unitary peripheral eigenmatrix exists]%
    \label{thm:local_symmetry_iff_spectral_radius}
    \lean{MPSTensor.localSymmetry_iff_spectralRadius_one}
    \leanok
    \uses{def:local_symmetry, def:twisted_transfer_map, def:injective,
          thm:twisted_transfer_spectral_radius_le_one}
    Let $A$ be injective, and assume $u u^\dagger = \Id$, $\E_A(\Id)=\Id$,
    $\Lambda$ is positive definite, $\tr(\Lambda)=1$, and
    $\E_A^\dagger(\Lambda)=\Lambda$.
    Then $u$ is a local symmetry if and only if there exist a unitary
    $V \in \MN{D}$ and a phase $\mu$ with $|\mu|=1$ such that
    \[
        \E_u(V) = \mu V.
    \]
    By Theorem~\ref{thm:twisted_transfer_spectral_radius_le_one}, this witness
    form is equivalent to the peripheral-spectrum condition $\rho(\E_u)=1$.
\end{theorem}

\begin{proof}
    \leanok
    \uses{thm:twisted_transfer_eigen_of_virtual_unitary,
        thm:twisted_transfer_modulus_one_gauge_phase,
        thm:virtual_unitary_of_twisted_gauge_phase,
        thm:boundary_state_invariant_of_virtual_unitary}
    The forward direction is Theorem~\ref{thm:twisted_transfer_eigen_of_virtual_unitary}.
    For the reverse direction, start from a unitary peripheral eigenmatrix.
    Theorem~\ref{thm:twisted_transfer_modulus_one_gauge_phase} turns it into a
    gauge-phase equivalence with the companion tensor; then
    Theorem~\ref{thm:virtual_unitary_of_twisted_gauge_phase} yields a virtual
    unitary and phase satisfying the local covariance relation, and
    Theorem~\ref{thm:boundary_state_invariant_of_virtual_unitary} gives the
    boundary-state invariance needed to assemble the local-symmetry data.
\end{proof}

\begin{theorem}[String order $\Leftrightarrow$ local symmetry]%
    \label{thm:string_order_iff_local_symmetry}
    \lean{MPSTensor.stringOrder_iff_localSymmetry}
    \leanok
    \uses{def:has_string_order, def:local_symmetry, def:injective}
    For an injective MPS, assume $u u^\dagger = \Id$,
    $\E_A(\Id)=\Id$, $\Lambda$ is positive definite,
    $\tr(\Lambda)=1$, and $\E_A^\dagger(\Lambda)=\Lambda$. Then string
    order for $u$ exists iff $u$ is a local symmetry.
\end{theorem}

\begin{proof}
    \leanok
    \uses{thm:string_order_implies_twisted_gauge_phase,
        thm:virtual_unitary_of_twisted_gauge_phase,
        thm:boundary_state_invariant_of_virtual_unitary,
        lem:has_string_order_of_local_symmetry}
    If string order exists, Theorem~\ref{thm:string_order_implies_twisted_gauge_phase}
    gives a gauge-phase equivalence with the companion tensor.
    Theorem~\ref{thm:virtual_unitary_of_twisted_gauge_phase} extracts a virtual
    unitary and phase from that equivalence, and
    Theorem~\ref{thm:boundary_state_invariant_of_virtual_unitary} shows that the
    same unitary preserves the boundary state, so one obtains local symmetry.
    Conversely, Lemma~\ref{lem:has_string_order_of_local_symmetry} turns local
    symmetry into a non-decaying boundary string-order sequence.
\end{proof}

\begin{theorem}[Virtual unitary from string order]%
    \label{thm:virtual_unitary_of_string_order}
    \lean{MPSTensor.virtualUnitary_of_stringOrder}
    \leanok
    \uses{def:has_string_order, def:injective}
    For an injective MPS, let $\Lambda \in \MN{D}$, and assume
    $u u^\dagger = \Id$ and $\E_A(\Id)=\Id$.
    If string order exists for $u$ with boundary parameter $\Lambda$, then
    there exists a virtual unitary $V$ and a unit-modulus scalar $\mu$ such that
    $\sum_j u_{ij} A^j = \mu\, V A^i V^\dagger$ for all~$i$.
    Diagrammatically this is the same local intertwining relation,
    but only up to an overall unit scalar on the virtual side:
    \begin{center}
        \TNCondCOne{u}{V}{i}
    \end{center}
\end{theorem}

\begin{proof}
    \leanok
    \uses{thm:string_order_implies_twisted_gauge_phase,
        thm:virtual_unitary_of_twisted_gauge_phase}
    String order first gives gauge-phase equivalence with the companion tensor by
    Theorem~\ref{thm:string_order_implies_twisted_gauge_phase}.
    Theorem~\ref{thm:virtual_unitary_of_twisted_gauge_phase} then normalizes the
    resulting intertwiner to a unitary and produces the phase $\mu$.
\end{proof}

%% ---------------------------------------------------------
\section{SPT Phase Classification}
%% ---------------------------------------------------------

Two injective symmetric tensors belong to the same symmetry-protected
topological phase when their virtual representation cocycles are
cohomologous~\cite{ChenGuWen2011}.

\begin{definition}[Same SPT phase]\label{def:is_same_spt_phase}
    \lean{MPSTensor.IsSameSPTPhase}
    \leanok
    \uses{def:cohomologous_to, def:projective_rep}
    Injective tensors $A$ and $B$, both symmetric under an on-site
    representation $U\colon G\to\GL_d(\C)$, are in the \emph{same SPT
        phase} if there exist virtual cocycles $\omega_A,\omega_B$ with
    corresponding projective representations that intertwine the
    respective twisted tensors and satisfy
    $\omega_A\sim\omega_B$ (cohomologous).
\end{definition}

\begin{theorem}[String order holds for injective symmetric MPS]%
    \label{thm:has_string_order_of_symmetric_injective}
    \lean{MPSTensor.hasStringOrder_of_symmetric_injective}
    \leanok
    \uses{def:has_string_order, thm:virtual_rep_injective}
    Let $A$ be an injective tensor, symmetric under a unitary
    representation $U$. Assume $\E_A(\Id)=\Id$, $\Lambda$ is positive
    definite, $\tr(\Lambda)=1$, and $\E_A^\dagger(\Lambda)=\Lambda$.
    Then string order exists for every group element $g$.
\end{theorem}

\begin{proof}\leanok
    \uses{thm:virtual_rep_injective,
        thm:twisted_transfer_modulus_one_gauge_phase,
        thm:virtual_unitary_of_twisted_gauge_phase,
        thm:boundary_state_invariant_of_virtual_unitary,
        lem:has_string_order_of_local_symmetry}
    The virtual representation theorem produces
    $\rho(g^{-1})\in\GL_D(\C)$ satisfying
    $\sum_j U(g)_{ij}\,A^j = \rho(g^{-1})\,A^i\,\rho(g^{-1})^{-1}$.
    A direct computation using $\E_A(\Id)=\Id$ shows that
    $\E_{U(g)}(\rho(g^{-1})) = \rho(g^{-1})$, so
    $\rho(g^{-1})$ is a nonzero eigenvector with eigenvalue~$1$.
    Theorem~\ref{thm:twisted_transfer_modulus_one_gauge_phase} upgrades this to
    gauge-phase equivalence with the companion tensor, and
    Theorem~\ref{thm:virtual_unitary_of_twisted_gauge_phase} extracts a virtual
    unitary and phase.
    By Theorem~\ref{thm:boundary_state_invariant_of_virtual_unitary}, this
    unitary preserves the boundary state, so
    Lemma~\ref{lem:has_string_order_of_local_symmetry} gives string order.
\end{proof}

\begin{theorem}[String order is an SPT invariant]%
    \label{thm:string_order_invariant_of_same_phase}
    \lean{MPSTensor.stringOrder_invariant_of_samePhase}
    \leanok
    \uses{def:is_same_spt_phase, thm:has_string_order_of_symmetric_injective}
    If $A$ and $B$ are injective, both on-site symmetric under a
    unitary representation $U$, and if $\E_A(\Id)=\Id$,
    $\E_B(\Id)=\Id$, $\Lambda_A$ and $\Lambda_B$ are positive
    definite, $\tr(\Lambda_A)=1$, $\tr(\Lambda_B)=1$,
    $\E_A^\dagger(\Lambda_A)=\Lambda_A$, and
    $\E_B^\dagger(\Lambda_B)=\Lambda_B$, then for every $g \in G$,
    string order exists for $A$ with twist $U(g)$ if and only if it
    exists for~$B$. In particular, since
    Definition~\ref{def:is_same_spt_phase} implies on-site symmetry
    for both tensors, this theorem applies to tensors in the same SPT
    phase.
\end{theorem}

\begin{proof}\leanok
    \uses{thm:has_string_order_of_symmetric_injective}
    Both directions follow from
    Theorem~\ref{thm:has_string_order_of_symmetric_injective}, which
    shows that string order holds universally for injective symmetric
    tensors with canonical normalisations.
\end{proof}
